\begin{abstract}
In January 2026 India's quick-commerce platforms withdrew their ``10-minute delivery'' claims after the Union Labour Ministry raised rider-safety concerns, and in August 2026 a Parliamentary Standing Committee recommended curbs on ``unrealistic speed targets'' without defining them. We combine a bibliometric review of operations research on last-mile and on-demand delivery \todo{N articles, 2015--2026} with an integrated optimisation model of a Bengaluru dark-store network built entirely from public data: WorldPop 2025 population, BBMP ward boundaries, OpenStreetMap retail sites and operating metrics from listed platforms' shareholder letters. A coverage-constrained facility-location model opens stores and allocates demand at minimum store and rider cost while guaranteeing a promise $T$ at a safe rider speed $v$; a transportation model plans replenishment, an assignment model dispatches riders at peak, and a PERT model estimates on-time probability. At \BaseV~km/h, the cost-optimal network has about \StoresTFifteen{} stores and reaches \ServedFifteenPct\% of residents for any promise of 15 minutes or more; with the available retail sites a 10-minute promise needs \StoresTTen{} stores and still leaves \UnreachTenPct\% unreachable. Serving everyone within 10 minutes requires \GridStorePremiumPct\% more stores and \GridCostPremiumPct\% higher cost per order than within 15 minutes. Rider speed, not the promise, drives cost: capping speed at 15 rather than 20~km/h raises cost per order by about \SpeedCapPremiumPct\%. Asking a 15-minute network to deliver in 10 minutes meets the promise for only \AssignStressTenOnTime\% of orders at a stressed store unless riders exceed safe speeds. We propose reachability within a safe ride radius as an operational definition of a ``realistic'' promise.

\medskip\noindent\textbf{Keywords:} quick commerce; dark stores; facility location; transportation problem; assignment problem; PERT; bibliometric analysis; gig-worker safety; India
\end{abstract}
